% Generated draft: requires mathematical and source review.
\documentclass[11pt]{article}
\usepackage[margin=25mm]{geometry}
\usepackage{fontspec}
\setmainfont{Latin Modern Roman}
\usepackage{amsmath,amssymb,mathtools,unicode-math}
\setmathfont{Latin Modern Math}
\usepackage{enumitem,needspace,xurl,hyperref}
\setlength{\parindent}{0pt}
\setlength{\parskip}{5pt}
\setlength{\emergencystretch}{3em}
\newcommand{\N}{\mathbb{N}}
\newcommand{\Z}{\mathbb{Z}}
\newcommand{\Q}{\mathbb{Q}}
\newcommand{\R}{\mathbb{R}}
\newcommand{\C}{\mathbb{C}}
\newcommand{\F}{\mathbb{F}}
\newcommand{\A}{\mathbb{A}}
\newcommand{\PP}{\mathbb{P}}
\newcommand{\E}{\mathbb{E}}
\newcommand{\eps}{\varepsilon}
\newcommand{\dd}{\,\mathrm{d}}
\begin{document}
\section*{UnsolvedMath ID: 700763}
\textbf{Source problem code:} catalogue entry
\subsection{Definitions and mathematical statement}
Write \(\N_{>0}=\{1,2,3,\ldots\}\). For subgroups \(U\leq H\trianglelefteq A\) of a finite group \(A\), put
\[
U^a=a^{-1}Ua\qquad(a\in A).
\]
The subgroup \(U\) is an \(A\)-covering subgroup of \(H\) if
\[
\bigcup_{a\in A}U^a=H.
\]
For finite groups, \([H:U]\) denotes the number of cosets of \(U\) in \(H\).

Does there exist a function \(f:\N_{>0}\to\N_{>0}\) such that, for every finite group \(A\), every normal subgroup \(H\trianglelefteq A\) with \(H<A\), and every proper subgroup \(U<H\),
\[
\bigcup_{a\in A}U^a=H
\quad\Longrightarrow\quad
[H:U]\leq f([A:H])?
\]
Thus the proposed bound depends only on \(n=[A:H]\), which satisfies \(n\geq2\), and is uniform over all such triples \((A,H,U)\).

The proposers report an affirmative answer when \(U\) is a maximal subgroup of \(H\).

The supplied \emph{Comments of 2025} report that the conjecture is proved when \(H\) acts innately transitively on the coset space \([H:U]\), citing M. Fusari, A. Previtali and P. Spiga; ``for \(n=3\) with \([H:U]\leq10\)'', citing L. Gogniat and P. Spiga; and when \(H=UL\), where \(L\) is a minimal \(A\)-invariant subgroup of \(H\), citing M. Fusari, S. Harper and P. Spiga. These are supplied literature claims, not independently verified conclusions. The wording of the \(n=3\) claim does not specify whether \([H:U]\leq10\) is a proved bound or an additional restriction. The final claim does not explicitly specify whether minimality excludes the trivial subgroup. These qualifications affect the reported special cases and do not alter the main question above.
\subsection{Short English statement}
Is the index of a proper \(A\)-covering subgroup \(U\) in a proper normal subgroup \(H\) of a finite group \(A\) bounded solely in terms of \([A:H]\)? The supplied registry labels the problem open and reports several affirmative special cases, with unresolved wording in two of those reports.
\subsection{Sources}
\begin{itemize}
\item Local UnsolvedMath registry, permanent ID \(700763\), exact source code \texttt{catalogue entry}, titled \emph{Kourovka Notebook Problem 11.71}. The supplied record labels the problem \emph{open}; this is a supplied status, not an independently verified assessment.
\item P. M. Neumann and C. E. Praeger, proposers, 1990. Source attribution: \emph{The Kourovka Notebook (Unsolved Problems in Group Theory)}, 11th issue (1990), Problem 11.71; supplied edition reference: 21st issue edition (2026), \texttt{arXiv:1401.0300v48}. The supplied collection description names E. I. Khukhro and V. D. Mazurov as editors.
\item M. Fusari, A. Previtali and P. Spiga, \emph{J. Group Theory} \textbf{27} (2024), 929--965. Supplied reference for the innately transitive case; no article title or link is supplied.
\item L. Gogniat and P. Spiga, \emph{Preprint}, 2025, \url{https://arxiv.org/abs/2502.01287}. Supplied reference for the statement ``for \(n=3\) with \([H:U]\leq10\)''.
\item M. Fusari, S. Harper and P. Spiga, \emph{Bull. Austral. Math. Soc.}, 2025, \url{https://doi.org/10.1017/S0004972725000176}. Supplied reference for the case \(H=UL\), with \(L\) a minimal \(A\)-invariant subgroup of \(H\).
\item The special-case assessments are supplied under \emph{Comments of 2025}. The registry record was updated on \texttt{2026-10-16T14:00:00Z}; this timestamp is registry metadata, not evidence of independent source verification. No separate literature-verification date is supplied.
\end{itemize}
\end{document}
